\chapter{A Bit That Lasts Twenty Seconds}\label{ch:majorana}

\begin{etymology}
\textbf{Topology} comes from the Greek \emph{tópos}, ``place'', and \emph{lógos},
``study''. It is the part of mathematics that studies what does not change when a
shape is stretched or bent without tearing: a coffee mug and a doughnut have one
hole each, and no small push can change that number.

\textbf{Parity} comes from the Latin \emph{paritas}, ``equality'', from \emph{par},
``equal''. In arithmetic it says whether a number is even or odd. Here it says whether
a fermion mode is empty or full.
\end{etymology}

\section*{Before this chapter}

This chapter is the story of a claim, and it does not stop to explain the physics
behind it. That is done in two earlier sections, and without them most of what follows
will be hard to follow:

\begin{itemize}
\item Section~\ref{sec:identical}, on identical particles: what bosons, fermions and
anyons are, and why exchanging two particles can apply a gate.
\item Section~\ref{sec:topological}, on topological qubits: the superconducting gap and
quasiparticles, the Majorana zero mode and its parity, the recipe for the wire, the
tetron, and how a parity is read through a quantum dot.
\end{itemize}

It also uses the $X$ and $Z$ of Chapter~\ref{ch:state} and the Gottesman--Knill theorem
of Chapter~\ref{ch:hardware}.

\bigskip

On 16 November 2026 Microsoft posted an eight-minute video called \emph{Majorana 2
Explained: The road to 1000x improvement}. In it, the people behind the company's
quantum chip say that its qubits became a thousand times more stable in one year,
that one operation now fits tens of millions of times inside the life of a qubit,
and that a useful machine is coming in 2029. Chapter~\ref{ch:hardware} explained how
the device is supposed to work, and ended its checklist with ``not shown''. This
chapter follows the bet from its start, in 1937, to the video, and then reads the
video against the paper behind it.

\section{A particle that is its own antiparticle}

\begin{history}{1937 \textperiodcentered{} Rome and Naples}
Ettore Majorana, a Sicilian physicist from Enrico Fermi's group in Rome, found a
version of Dirac's equation whose solutions are their own antiparticles: a particle
and its mirror image are the same object. He had the result for years and did not
publish it. He brought it out in \emph{Il Nuovo Cimento} in 1937, when colleagues
convinced him to compete for a professorship, and was appointed at Naples for his
reputation alone. On 25 March 1938 he took a boat from Palermo to Naples and was
never seen again. His body was never found.
\end{history}

Whether any elementary particle is a Majorana fermion is still open; the neutrino is
the candidate, and experiments are still looking. The story of this chapter is about
something else with the same name: the half-electron, half-hole excitations at the ends
of a superconducting wire, which are their own antiparticles in the sense of
Section~\ref{sec:topological}.

\section{Kitaev's idea}

\begin{history}{1997 \textperiodcentered{} Landau Institute, Moscow}
Alexei Kitaev, the same Kitaev of phase estimation (Chapter~\ref{ch:fourier}), and of the
surface codes of Chapter~\ref{ch:qec}, posted \emph{Fault-tolerant quantum computation
by anyons} in July 1997. His proposal: store qubits in non-abelian anyons, whose shared
state depends only on how they were wound around each other, the topology of their
paths. Moving one around another applies a gate, and the gate is exact as long as the
winding is right, however shaky the hand that did it.
\end{history}

In October 2000, by then also affiliated with Microsoft Research, Kitaev wrote down the
wire model of Section~\ref{sec:topological}, with one Majorana zero mode at each end. The
promise attracted Microsoft early. By 2006 the company had a group at the University of
California, Santa Barbara, led by the Fields medalist Michael Freedman, devoted to
topological quantum computing. When the people in the 2026 video say ``we've worked on
topological quantum computing for 20 years'', this is the start they mean.

\section{The peak}

In 2010 two groups published the recipe of Section~\ref{sec:topological}: a semiconductor
nanowire, a superconducting coat and a magnetic field. Experimenters started right away.

\begin{history}{2012 \textperiodcentered{} Delft}
Vincent Mourik, Leo Kouwenhoven and colleagues followed the recipe with indium antimonide
wires and a superconducting contact. They sent current into the end of the wire and saw a
peak of conductance at exactly zero voltage, staying put as they changed the field and the
gate voltages. A Majorana mode at the end of a wire should produce exactly that peak. The
paper appeared in \emph{Science} and started a race.
\end{history}

The trouble is that the Andreev bound states of Section~\ref{sec:topological} produce the
same peak with no topology at all. The following decade was spent trying to tell the two
apart, and most of what follows in this chapter is that same argument, at higher and higher
stakes.

\section{The retraction}

\begin{history}{2018--2021 \textperiodcentered{} Delft}
In 2018 Kouwenhoven's group, partly funded by Microsoft, published \emph{Quantized Majorana
conductance} in \emph{Nature}: the zero-voltage peak had reached the exact height theory
predicts for a Majorana mode, $2e^2/h$. Sergey Frolov and Vincent Mourik, who had worked on
the 2012 paper, asked for the raw data and found that data contradicting the claim had been
left out. On 8 March 2021 the authors retracted the paper. After a new calibration the
plateaus sat 8\% higher than the quantized value, and the authors wrote that they could no
longer claim the observation. An investigation commissioned by Delft found no fabrication,
but found some selection in what had been published.
\end{history}

The retraction did not end the program. It changed its strategy. Microsoft moved most of
the work in house and replaced the single zero-voltage peak with a battery of measurements.

\section{Majorana 1}

In 2022 Microsoft posted the \emph{topological gap protocol}, published in \emph{Physical
Review B} in 2023: a sequence of transport measurements on both ends of a wire that, the
company says, certifies that a region of magnetic field and gate voltage is in the
topological phase. Only devices that pass the protocol move on to the qubit experiments.

On 19 February 2025 Microsoft announced \emph{Majorana 1}, a chip with eight qubits and,
according to the company, a path to a million. The same day \emph{Nature} published a
paper with the measurement behind it: a single-shot measurement of the parity of a wire,
read through a quantum dot whose capacitance changes with the parity. The paper did not
claim the announcement. The editors added a note to the peer review file saying that the
results ``do not represent evidence for the presence of Majorana zero modes in the reported
devices'', and that the paper was published for its device architecture. The press release
said the topological qubit had arrived. The paper said a parity had been read.

Henry Legg, at the University of St Andrews, posted two comments within a month. In the
first he argued that the topological gap protocol has no consistent definition of
``gap'' or of ``topological'', and that its verdict changes with the ranges of field,
voltage and resolution chosen for the analysis. In the second he showed that the same
protocol can classify the regions where the 2025 parity readout was done as either gapped
or gapless.

A qubit needs both loops of the tetron in Figure~\ref{fig:tetron}, because a bit that can
only be read in $Z$ is a classical bit. In July 2025 Microsoft posted the first measurement
of both on one device. The $Z$ parity switched about every $12.4$~ms. The $X$ parity
switched about every $14.5\,\mu$s, almost a thousand times faster, and its readout was wrong
16\% of the time.

\section{Majorana 2}\label{sec:majorana2}

On 2 June 2026 Microsoft announced \emph{Majorana 2} and posted the paper behind it, with 165
authors: \emph{20 Second Parity Lifetime in an InAs--Pb Tetron Device}. Three things changed
at once. The superconductor went from aluminum to lead, whose gap is about eight times larger.
The semiconductor went from indium arsenide to a stack of indium arsenide and indium arsenide
antimonide. And the way the device is tuned changed to a radio-frequency method precise to
about a microelectronvolt. The topological gap that the company reports rose from about 30 to
about $70\,\mu$eV.

The headline number is the $Z$ parity: it now switches on average every $22 \pm 1$~s, some
dwell times reach minutes, against 1 to 12~ms in the aluminum devices. That is where the
``1000x'' comes from; against the 12~ms device the ratio is 1833, and against the 1~ms device
it is 22\,000. The paper is a preprint. In November 2026 it had not been through peer review.

\begin{reflection}{A gap twice as large gives a lifetime a thousand times longer?}
The video says so directly: the material improved ``by a multiple of a few, two to three'',
and that turned into a factor of a thousand. It sounds like an exaggeration, and it is the
most ordinary part of the story.

The $Z$ parity flips when something supplies enough energy to cross the gap. Thermal
processes of that kind have a rate proportional to $e^{-\Delta/k_B T}$, with $\Delta$ the gap
and $T$ the temperature of whatever is causing the flips. Going from $\Delta = 30$ to
$70\,\mu$eV multiplies the lifetime by $e^{40\,\mu\mathrm{eV}/k_B T}$. For that factor to be
1000, $40\,\mu\mathrm{eV}/k_B T = \ln 1000 \approx 6.9$, which gives $k_B T \approx
5.8\,\mu$eV, or $T \approx 67$~mK. A dilution refrigerator sits at about 15~mK, and the stray
excitations in the device are usually hotter than the refrigerator. Sixty-seven millikelvin is
a reasonable number for them.

So a factor of two in an exponent gives a factor of a thousand outside it. That is the same
arithmetic that makes the Boltzmann factor in Chapter~\ref{ch:hardware} drop from a coin toss
at 300~K to almost nothing at 15~mK. The thousand is real. What it measures is the next
question.
\end{reflection}

Three weeks later, on 24 June 2026, \emph{Nature} published Legg's critique of the 2025 paper
as a \emph{Matters Arising}, with Microsoft's reply in the same issue. Legg had gone through
the code of the topological gap protocol, which Microsoft had made public. A plotting filter
was hard-coded to show only the largest region that passed the protocol, hiding a second one;
and the code flipped the bias voltage by position in an array rather than by its physical
value. With those fixed, he got different results from the same measurements, and he wrote
that the coding errors had led to incorrect statements to the peer reviewers. Microsoft
called it ``a minor off-by-one-pixel bug'', said that the protocol is only a tuning step with
no role in interpreting the qubit measurements, and said that Legg offered no alternative
physical model that reproduces their capacitance signal. On Majorana 2, Legg told \emph{The
Register} that it ``is not proven to even be a single qubit''.

The theorists answered too. In August 2026 Binayyak Roy, Jay Sau and Sumanta Tewari showed
that a short segment of the wire in the ordinary, non-topological state can produce the same
flux-dependent capacitance signal that Microsoft reads as the signature of Majorana modes, and
proposed a test to tell them apart. Jay Sau is a coauthor of the 2010 recipe. The people who
invented the platform are among the ones asking for more evidence.

\section{What the video says}

\begin{question}{Is a twenty-second qubit better than a millisecond transmon?}
It is not a twenty-second qubit. It is a twenty-second bit. The number measures how long the
$Z$ parity stays put, which is how long a classical 0 or 1 survives without an outside
electron jumping in. A qubit also has to keep its $X$ value, which is where the phase lives,
and Majorana 2 does not report $X$. The transmon's millisecond in Chapter~\ref{ch:hardware} is
measured on a working qubit with gates. The two numbers are not about the same thing.
\end{question}

The video does its arithmetic in a single sentence: with an operation every microsecond and a
qubit that lives a millisecond, there is one error every thousand operations; with tens of
seconds, ``you can do tens of millions of operations before an error''. With the measured
numbers and the same yardstick of one microsecond per operation:

\begin{center}
\small
\renewcommand{\arraystretch}{1.2}
\begin{tabular}{llrr}
\hline
device & parity & lifetime & operations per lifetime \\
\hline
aluminum tetron (2025) & $Z$ & 12.4 ms & 12\,400 \\
aluminum tetron (2025) & $X$ & 14.5\,$\mu$s & about 15 \\
lead tetron (2026) & $Z$ & 22 s & $2 \times 10^7$ \\
lead tetron (2026) & $X$ & not reported & -- \\
\hline
\end{tabular}
\end{center}

The tens of millions are real for $Z$. The last time anyone published $X$, the same sentence
would have said fifteen. Legg put it in one question: why do they not report the $X$
measurement? The lead may well have improved $X$ too; the video speaks as if it had, and the
paper does not show it. Saying ``qubits 1000 times more reliable'' about a measurement of one
of the two values a qubit needs is marketing.

The rest of the video is the usual. Quantum computers will feed ``the exact, ground-truth
answer'' into AI models so that Copilot can ``create materials and molecules''; they will find
superconductors that work at room temperature. Part~\ref{part:problems} is about the problems
whose cost is actually known, and these are not among them. The one honest line is about
skepticism: ``there's many things that are proprietary, that we can't share all the details''.
Legg found the bugs because the code of the protocol was public. A claim this size, from a
program with a retraction behind it, needs more data, not less.

\section{If it works}

Suppose the next paper shows a long $X$ lifetime, measurements that are right almost every time,
and a test that rules out the ordinary explanations. There would be a topological qubit, and it
would still not be a quantum computer. As Section~\ref{sec:topological} explained, measuring and
exchanging Majorana modes gives Clifford gates and nothing else, and the $T$ gate has to come from
magic states, prepared outside the protection and purified at the cost that Part~\ref{part:qec}
counts. Topological protection buys criterion 3 and a smaller overhead per logical qubit. It does
not buy universality.

Then there is the date. Microsoft's blog says the company ``cut our timeline in half'' and aims for
a scalable quantum computer by 2029. What the United States government has done is narrower. In
DARPA's Quantum Benchmarking Initiative, Microsoft is in Stage C, the final one, which it entered
through an earlier pilot program. Stage C is an independent check of whether the proposed machine
can be built as designed. On 7 October 2026 DARPA added Atom Computing, Diraq, IBM and IonQ to it,
alongside Microsoft and PsiQuantum, and the program's own target for a utility-scale machine is
2033. Being in Stage C means the design is being audited. It does not mean the qubit works.

The row for topological qubits in the table of Chapter~\ref{ch:hardware} still says ``not shown''.
After Majorana 2 it is still the honest entry. What changed in 2026 is the material: a lead
superconductor and a gap more than twice as large gave a bit that holds for twenty seconds, and
that is a real result. The qubit is still waiting for its $X$.
